\chapter{State of the Art and Technology Stack}

\section{Smart Home Platforms: Current Landscape}
Smart-home systems generally follow one of three patterns:
\begin{enumerate}[leftmargin=1.2cm]
  \item Vendor-locked ecosystems focused on proprietary devices.
  \item Gateway-centric open-source stacks with limited user governance.
  \item Cloud-native IoT services with strong scaling but high integration overhead.
\end{enumerate}

The SmartHome platform targets a hybrid engineering path: flexible protocol-level integration (MQTT), application-level governance (REST + role model), and web-native UX (React).

\section{Protocol and Interaction Model}
\subsection{MQTT for Device Messaging}
MQTT is lightweight, topic-based, and suitable for constrained devices and low-bandwidth communication. The project uses MQTT for sensor telemetry and device command delivery.

\subsection{HTTP/REST for Business Operations}
REST endpoints handle account management, CRUD workflows, permissions, and administrative functions where request-response semantics are natural.

\subsection{Socket.IO for User-facing Real-time Events}
Socket.IO complements REST by delivering asynchronous updates such as approval notifications, gas alerts, and dashboard state changes.

\section{Technology Selection Rationale}
\subsection{Backend: Node.js + Express}
Node.js provides asynchronous I/O efficiency for mixed workloads (HTTP, sockets, MQTT events). Express offers low ceremony and direct control over middleware and route composition.

\subsection{Database: MongoDB + Mongoose}
MongoDB supports document flexibility for evolving IoT payload fields and role-specific user metadata. Mongoose adds schema constraints and lifecycle hooks (e.g., password hashing).

\subsection{MQTT Broker: Aedes}
Aedes allows embedding broker behavior directly inside backend runtime, reducing integration complexity in prototyping and early production stages.

\subsection{Frontend: React + Vite}
React enables component-level decomposition and route guards. Vite ensures fast development iteration and efficient modern build tooling.

\subsection{Realtime: Socket.IO}
Socket.IO handles transport fallback and namespace/room abstractions, simplifying per-house and per-role event fan-out.

\section{Comparative Positioning}
\begin{table}[H]
\centering
\caption{Comparative perspective of architectural choices}
\begin{tabular}{p{4cm}p{4cm}p{6cm}}
\toprule
Dimension & Typical Alternative & Adopted Choice and Advantage \\
\midrule
Device Messaging & HTTP polling from devices & MQTT topics: lower overhead, event-native pattern \\
Realtime UI & Frequent client polling & Socket push model: lower latency and less load \\
Data Model & Strict relational schema & Flexible documents for IoT and user metadata \\
Auth Model & Single-role simple auth & Multi-role with permission workflow and onboarding \\
\bottomrule
\end{tabular}
\end{table}

\section{Engineering Standards and Principles}
The project follows practical engineering principles:
\begin{itemize}[leftmargin=1.2cm]
  \item Separation of concerns across route, model, middleware, and realtime services.
  \item Explicit role checks for privileged operations.
  \item Deterministic state transitions for safety-critical workflows.
  \item Event-driven synchronization between data persistence and UI state.
\end{itemize}

\section{Toolchain and Dependencies}
\subsection{Backend Core Libraries}
Important dependencies include Express, Mongoose, Aedes, Socket.IO server, JWT, bcrypt, and OAuth strategy packages.

\subsection{Frontend Core Libraries}
Important dependencies include React, React Router, Axios, Socket.IO client, MQTT client, and chart/animation utilities.

\subsection{Development Tooling}
Nodemon and Vite support development cycles; ESLint supports code quality baselines.

\section{Known Trade-offs}
\begin{itemize}[leftmargin=1.2cm]
  \item Embedded broker simplifies architecture but couples API and broker runtime.
  \item Flexible schemas ease evolution but require strict validation discipline.
  \item Socket-heavy UX improves responsiveness but demands robust auth handling and room hygiene.
\end{itemize}

\section{Chapter Synthesis}
The selected stack is coherent with project constraints: quick iteration, protocol-level interoperability, and real-time operational responsiveness. Subsequent chapters translate these choices into formal architecture and implementation.
